\chapter{Energy Sources and RF Harvesting}
\label{ch:harvest}

\begin{defbox}[In this chapter you will learn]
\begin{itemize}
 \item which ambient energy sources can power a node, and roughly how much each delivers;
 \item how a rectenna turns RF into DC and why its efficiency depends strongly on input power;
 \item Amber's harvesting model: how received power becomes an \emph{equivalent source voltage} behind a series resistor, and what that implies about how much energy reaches the capacitor;
 \item the three \emph{power modes}---wireless power transfer, external, hybrid---and the difference between combining them by \emph{max} and by \emph{sum};
 \item how an external energy trace is replayed, and a surprise about the shipped trace.
\end{itemize}
\end{defbox}

\section{Where can the energy come from?}
\label{sec:sources}

``Ambient'' in Ambient IoT refers to the environment as the supplier of energy
and, sometimes, of the carrier. The main candidates differ by orders of magnitude.

\begin{center}
\small
\renewcommand{\arraystretch}{1.15}
\begin{tabularx}{\linewidth}{@{}>{\sffamily\bfseries}p{3.4cm}>{\raggedright}p{3.8cm}X@{}}
\toprule
Source & Typical power density & Remarks\\
\midrule
Outdoor sunlight & $\sim\SI{10}{mW/cm^2}$ and more & Very high, intermittent (day/night, shade)\\
Indoor light & $\sim\SI{10}{}$--$\SI{100}{\micro W/cm^2}$ & Moderate; depends on lighting\\
Thermal gradients & tens of $\si{\micro W/cm^2}$ & Needs a temperature difference\\
Vibration & $\sim\SI{100}{\micro W/cm^3}$ & Machinery, vehicles\\
Ambient RF & $\ll\SI{1}{\micro W/cm^2}$ & Very low except near a transmitter\\
\emph{Dedicated} RF (carrier) & depends on distance & Controlled; what Amber's base stations provide\\
\bottomrule
\end{tabularx}
\end{center}
\noindent\emph{Order-of-magnitude values commonly quoted in the energy-harvesting literature; actual numbers vary widely with conditions.}

Ambient RF alone is too weak to be dependable, so 3GPP Ambient IoT relies on
the network itself to \emph{beam} energy at the devices: the base station's carrier
is both the power source and the illumination for backscatter. A node can
additionally use a second source---typically light---so that it can wake up even
where the RF field is weak. Amber lets you model exactly that combination.

\section{The RF harvesting chain}
\label{sec:rectenna}

A \emph{rectenna} (rectifying antenna) consists of
\begin{enumerate}
 \item the antenna, which captures the RF power $P_{\mathrm{dl}}$;
 \item an impedance-matching network;
 \item a rectifier---diodes or transistors configured as a voltage multiplier---that converts the oscillating RF voltage into DC;
 \item a smoothing capacitor and a power-management stage that feed the storage element.
\end{enumerate}
The \emph{RF-to-DC conversion efficiency} $\eta_{\mathrm{rf\text{-}dc}}=P_{\mathrm{dc}}/P_{\mathrm{rf}}$
depends strongly on the input power, because a diode barely conducts until the
RF voltage exceeds its threshold. At very low input power (below roughly
$\SI{-20}{dBm}$ for commonly used devices) the efficiency collapses, and many
rectifiers cannot output anything usable at all---the \emph{cold-start}
problem. At moderate power the efficiency peaks (tens of percent) and at high
power the diode breaks down. A rectenna therefore has a \emph{minimum input power}
below which the harvested energy is effectively zero, however long you wait.
This is why, in a real deployment, the \emph{forward link} often limits range.

\begin{reality}[Amber has no rectifier curve]
In Amber every watt of received RF power reaches the harvester: effectively
$\eta_{\mathrm{rf\text{-}dc}}=1$ at all input levels, with no turn-on threshold,
no saturation and no cold start. The code comment simply converts received
power to DC power. If you need a realistic harvester, the hook is the single
conversion in \codew{CoverageMap.calculate_node_power} (Exercise~\ref{ex:rectifier}).
\end{reality}

\section{Amber's harvesting model}
\label{sec:amberharvest}

\subsection{From power to an equivalent source}

The capacitor model of Chapter~\ref{ch:capacitor} is driven by a
\emph{voltage} source in series with a resistor $\Rs$. The controller therefore
turns the harvested power $P_h$ (a number in dBm stored in \codew{node.harvesting_power_dbm})
into an equivalent open-circuit voltage by
\begin{equation}
 \Vs=\sqrt{P_h\,\Rs},\qquad\text{i.e. } P_h=\frac{\Vs^{2}}{\Rs}.
\label{eq:vsrc}
\end{equation}
This is the usual $P=V^2/R$ read backwards. Physically it says: ``a source of
internal resistance $\Rs$ and open-circuit voltage $\Vs$ would \emph{dissipate} $P_h$
if short-circuited into a resistance equal to its own''. It is a modelling device, not a
circuit that exists on a board; but its consequences are physical and worth
understanding.

\begin{intuition}[How much of $P_h$ actually reaches the capacitor?]
When the capacitor sits at voltage $V$, the current is $(\Vs-V)/\Rs$ and the power
\emph{delivered into the capacitor} is $V(\Vs-V)/\Rs$. With $x=V/\Vs$,
\[
 P_{\mathrm{into\ C}}(x)=P_h\,x(1-x),
\]
which is a parabola with maximum $P_h/4$ at $x=\tfrac12$ (the familiar
\emph{maximum-power-transfer} theorem). So the \emph{effective} harvesting efficiency of
the model is at most $25\%$ and falls to zero as the capacitor approaches $\Vs$.
For Node~0 at $V=\SI{1.5}{V}$, $x=0.60$ and the capacitor receives
$0.24\,P_h\approx\SI{0.30}{mW}$ of the $\SI{1.24}{mW}$ nominally harvested; the
rest is dissipated in $\Rs$. The charge-only update of the simulator reproduces this behaviour exactly
(Chapter~\ref{ch:capacitor}).
\end{intuition}

\subsection{Three power modes}

The harvest power of each node is recomputed from the pre-computed downlink power
by \codew{CoverageMap.calculate_node_power}, according to \codew{node_energy_mode}.
Let $P_{\mathrm{dl}}$ be the best downlink power (in watts) at the node, $\eta$ its
\code{efficiency}, $P_{\mathrm{ext}}$ the instantaneous power of the external source,
and $\mathbb{1}_{\mathrm{tx}}$ be $1$ while the node is in its \emph{transmitting} state and $0$ otherwise.

\begin{center}
\small
\renewcommand{\arraystretch}{1.25}
\begin{tabularx}{\linewidth}{@{}>{\ttfamily}l>{\raggedright}X>{\raggedright\arraybackslash}X@{}}
\toprule
\normalfont\bfseries Mode & \normalfont\bfseries Harvested power $P_h$ & \normalfont\bfseries Backscatter power\\
\midrule
"wpt" & $P_{\mathrm{dl}}\,\big(1-\eta\,\mathbb{1}_{\mathrm{tx}}\big)$ & $\eta P_{\mathrm{dl}}$\\
"external" & $P_{\mathrm{ext}}$ (RF is \emph{not} harvested) & $\eta P_{\mathrm{dl}}$\\
"hybrid" & $\max\big(P_{\mathrm{wpt}},P_{\mathrm{ext}}\big)$ if \codew{combine_mode="max"} (default), else $P_{\mathrm{wpt}}+P_{\mathrm{ext}}$ & $\eta P_{\mathrm{dl}}$\\
\bottomrule
\end{tabularx}
\end{center}
where $P_{\mathrm{wpt}}=P_{\mathrm{dl}}(1-\eta\,\mathbb{1}_{\mathrm{tx}})$. Three points deserve emphasis.
\begin{enumerate}
 \item The backscatter power is $\eta P_{\mathrm{dl}}$ in \emph{every} mode: a node powered by solar light still
   needs the base station's carrier to reflect.
 \item ``Hybrid'' does \emph{not} mean ``sum''. With the default \codew{combine_mode="max"}
   the node simply uses whichever source is stronger at that instant.
 \item The base station carriers are not summed either: the node harvests the
   \emph{best} sector's power only (\code{best_pw_w}), a conservative simplification when several
   carriers overlap.
\end{enumerate}

\begin{example}[Why the hybrid node of the demos lives on sunlight]
In the broadcast demos Node~2, $\SI{100}{m}$ from the base station, receives
$P_{\mathrm{dl}}=\SI{-10.57}{dBm}=\SI{87.7}{\micro W}$, giving $\Vs=\sqrt{87.7\times10^{-6}\cdot5000}=\SI{0.662}{V}$ from RF. The steady-state voltage
(Chapter~\ref{ch:capacitor}) is $\Vs\Rl/(\Rl+\Rs)=\SI{0.63}{V}\ll\Vhigh$, so from RF alone Node~2 would
\emph{never} start. The solar trace has mean power $\SI{1.43}{mW}$ ($\Vs\approx\SI{2.66}{V}$),
$16\times$ larger. Because \codew{combine_mode="max"} picks the larger term, \emph{the RF contribution is
irrelevant}; Node~2's harvesting is solar-only, which is why it activates fastest
and sends the most packets (Figure~\ref{fig:sawtooth}). With \codew{combine_mode="sum"} the two would add in power:
$\sqrt{(87.7+1433)\,\si{\micro W}\cdot5000}=\SI{2.76}{V}$.
\end{example}

\section{The external source}
\label{sec:envsource}

\codew{EnvEnergySource} is a SimPy actor that \emph{replays a recorded trace}.
It reads an Excel file, takes one column (default \code{V_IM}) of voltages, and at
every simulation millisecond sets
\[
 P_{\mathrm{ext}}(t_k)=\frac{V_k^{2}}{R},\qquad R=\code{resistance}\ (\SI{1}{\ohm}\ \text{default; the demos pass }R_s=\SI{5}{k\ohm}).
\]
After the last row it loops. The object exposes the current value as
\code{ext_power}; the coverage map reads it through the callback
\codew{node_ext_power_fn}, which in the demos is
\codew{lambda node: energy_source.ext_power}. Because the callback ignores its
argument, \emph{all nodes see the same external power at the same instant}: there is no
spatial variation (shading) of the solar source.

When \code{resistance} equals the controller's $\Rs$, the round trip
$V\to P=V^2/R\to\sqrt{P\Rs}$ returns the recorded voltage unchanged. In the demos the trace's \code{V_IM}
column is thus directly the equivalent source voltage of the solar harvester.

\begin{reality}[The ``300-second'' trace is three seconds long]
The file \file{threehundred_seconds_stable.xlsx} contains $2999$ rows with a time
step of $\SI{1}{ms}$ (the last time stamp is $\SI{2.998}{s}$). Since
\codew{EnvEnergySource} consumes one row per millisecond and loops, the
trace is \emph{three seconds long}, repeated. Its statistics are: $V_{IM}$ between
$1.78$ and $\SI{3.01}{V}$, mean $\SI{2.66}{V}$, standard deviation $\SI{0.25}{V}$, i.e.\ a
nearly constant $\approx\SI{1.4}{mW}$ source ($\SI{1.56}{dBm}$). For simulations of tens of
seconds the periodicity of the trace could show up as spurious structure.
\end{reality}

\section{When is harvesting updated?}

Harvest power is recomputed \emph{every millisecond for every node} inside the
controller loop (Chapter~\ref{ch:controller}). For a pure WPT node this matters
only when its state toggles into or out of \emph{transmitting} (the $1-\eta$
factor); for a node with an external source it also tracks the trace. The comment
in the controller source explains the economy: if no coverage map or downlink
results are supplied, ``wpt is static, no need to recalculate every ms''.

\begin{codemap}
\small
\begin{tabularx}{\linewidth}{@{}>{\ttfamily}p{5.0cm}X@{}}
\toprule
\normalfont\bfseries Concept & \normalfont\bfseries Code\\
\midrule
Harvest/backscatter bookkeeping & \codew{CoverageMap.calculate_node_power} (\file{propagation.py})\\
Hybrid combination & \codew{combine_mode} (\code{"max"} default)\\
Power to source voltage, Eq.~\eqref{eq:vsrc} & \codew{Controller.run}: \codew{(power_w * R_series) ** 0.5}\\
Solar trace & \codew{EnvEnergySource} (\file{energy.py})\\
Callback & \codew{CoverageMap.node_ext_power_fn}\\
\bottomrule
\end{tabularx}
\end{codemap}

\section{Chapter summary}
\begin{itemize}
 \item Sunlight dwarfs ambient RF by orders of magnitude; 3GPP therefore beams a dedicated carrier at the devices.
 \item Real rectennas have a minimum input power and a power-dependent efficiency; Amber omits both.
 \item Amber converts harvested power to a Th\'evenin-style source $\Vs=\sqrt{P_h\Rs}$; at most $P_h/4$ ever reaches the capacitor.
 \item Modes: wpt, external, hybrid (max or sum). Backscatter always uses the base-station carrier.
 \item The shipped trace is $3$~s long, shared by all nodes, and effectively a $\SI{2.7}{V}$ source.
\end{itemize}

\section*{Exercises}
\begin{exercise}
The solar trace has mean voltage $\SI{2.66}{V}$ and $R=\SI{5}{k\ohm}$. Compute the external power in mW and dBm. (Careful: use the mean of $V^2$ or of $V$? How big is the difference given $\sigma=\SI{0.25}{V}$?)
\end{exercise}
\begin{exercise}
In hybrid mode, $P_{\mathrm{wpt}}=\SI{1}{mW}$ and $P_{\mathrm{ext}}=\SI{1.4}{mW}$. Compute the equivalent source voltage ($\Rs=\SI{5}{k\ohm}$) for \codew{combine_mode="max"} and for \code{"sum"}.
\end{exercise}
\begin{exercise}
A WPT node with $\eta=0.7$ receives $\SI{1}{mW}$. By how many dB does its harvested power drop while it is transmitting? What if $\eta=0.2$?
\end{exercise}
\begin{exercise}
For which received downlink power $P_{\mathrm{dl}}$ can a WPT-only node eventually reach $\Vhigh=\SI{1.7}{V}$ with $\Rs=\SI{5}{k\ohm}$ and $\Rl=\SI{100}{k\ohm}$? Express your answer in $\si{\micro W}$ and dBm. (Hint: you need the steady-state formula of Chapter~\ref{ch:capacitor}.)
\end{exercise}
\begin{exercise}\label{ex:rectifier}
Suppose you want to include an RF-to-DC efficiency $\eta_{\mathrm{rf\text{-}dc}}(P_{\mathrm{in}})$ that is zero below $\SI{-20}{dBm}$ and $0.3$ above. Where in the pipeline $P_{\mathrm{dl}}\to\Vs$ would you insert it, and how would this change the paper's scenario of Section~\ref{sec:twoway}?
\end{exercise}
